\documentclass[12pt]{article}

\usepackage[a4paper, total={6in, 8in}]{geometry}
\usepackage{textcomp}

\begin{document}
\setlength{\parindent}{0pt}

\def \t {\textrightarrow}
\def \v {\vspace{1em}}
\def \bi {\begin{itemize}}
\def \ei {\end{itemize}}
\def \s[#1] {\section*{#1}}
\def \ss[#1] {\subsection*{#1}}
\def \sss[#1] {\subsubsection*{#1}}

\s[Tacito]
Parla della germania e delle abitudini di una popolazione che c'entra con Pier Paolo Pasolini \\
Sono rozzi, una popolazione viscerale con delle regole precisissime \t in questa visceralita aveva un anima pura, che i romani non avevano piu \\
Roma era corrotto, teatro di delitti, fratricidi, matricidi etc. \t roma attenta ai disvalori \t le corti degli imperatori erano una monnezza, con sporcizia umana \\
Queste popolazioni sono invece pure da questo punto di vista \\
Qua c'è somiglianza con "I ragazzi di vita" di Pasolini \t parla di emarginati dalla societa, poveri ragazzi senza famiglia e senza istruzioni \t hanno vaolri viscerali, ma sono pronti ad accogliere come vasi \\
Tacito nasce nel 55 d.C. \t nasce a terni, ma andra anche in germania dove conosce queste popolazioni

\ss[Agricola]
Studia a roma e sposa la figlia di Giulio Agricola \t Tacito chiama "Agricola" un trattato celebrativo/encomiastico, circa la figura del suocero \\
"De Vita et de Moribus Iulii Agricolae" \t opera in cui esalta le qualita del suocero \t ma non è una celebrazione della morte, ma un incentivo alla classe dirigente del futuro \t fa parte della linea intellettuale-potere \\
L'opera mette in risalto tutte le virtu di Agricola \t molti hanno accusato Tacito di essere stato raccomandato dal suocero \\
Qui si coglie il suo essere eccezzionale \t ritratto come intellettuale e non come politico \\
\textbf{Se la germania è un opera etnogeografica e geoetnografica, Agricola è un opera encomistico/celebrativa}

\v

Va in germania \t dove le persone si lavavano insieme ed erano mezzi nudi \t con regole rigidissime anche con le donne \\
Ma erano genuine, ed erano autentici e spontanei \t tutto cio che manca a roma \t \textbf{tutto si riconduce alla visione di roma che cessa di essere romanocentrica, e diviene barbarocentrica} \t roma non è + la roma di augusto \\
Roma capisce che c'è altro \t si perde visione romano-centrica, ora etero-centrica \t questo generera grande fattore storico \t in Virgilio, per esempio, si aveva gia Didone che non era nata a Roma \t e gia questo spotava l'asse sull'altro da noi \\
Il fatto storico che si aggangia sono i regni romano-barbarici \t se roma è il punto di partenza per il sistema giuridico \\
Termini chiave: caput mundi, canto 6 dell eneide, regno barbaro-romano, barbarocentrico
\end{document}
